% Design note based on the discussion on 2026-10-15; not an experiment report.
% Section fragment: include inside an existing document with
% % Design note based on the discussion on 2026-10-15; not an experiment report.
% Section fragment: include inside an existing document with
% % Design note based on the discussion on 2026-10-15; not an experiment report.
% Section fragment: include inside an existing document with
% % Design note based on the discussion on 2026-10-15; not an experiment report.
% Section fragment: include inside an existing document with
% \input{sampling-training-evaluation-workflows}
% Needs Chinese-capable fonts, tikz, and \usetikzlibrary{arrows.meta}.
% The existing framework.tex already provides these prerequisites.

\section{比較兩階段 sampling 時的 training 與 evaluation 切分}
\label{sec:sampling-training-evaluation-workflows}

本節記錄預定的實驗設計，\allowbreak 不代表已經實作或執行。
這裡的 training records 是提供給 LLM 生成 description 的範例，\allowbreak
不會訓練或更新模型權重。
為避免混淆，\allowbreak 下文區分「training 候選池」與「最後送進 LLM 的 generation list」。

\begin{itemize}
  \item \textbf{第一次 sampling：}collector 在 proof 執行時選取並保留候選 records。
  \item \textbf{中間處理：}依既定流程合併候選、exact dedup 與 similarity check。
  \item \textbf{第二次 sampling：}從候選池依分組配額與遞補規則，\allowbreak
        選出最後用來生成 description 的 records。
\end{itemize}

\subsection{情況一：比較第一次 sampling 方法}
\label{subsec:holdout-first-sampling}

A、B 從相同 proof 與 source commit 開始，\allowbreak
使用不同的第一次 sampling 方法。
兩邊使用相同候選容量，\allowbreak 並固定其餘 collector 設定、dedup 與 similarity check 規則。
因此，中間處理後的候選 lists 可以不同，\allowbreak 這正是要保留的實驗差異。

\begin{center}
\begin{tikzpicture}[
  >=Latex,
  flow/.style={->, thick},
  box/.style={draw, rounded corners, align=center, text width=5.2cm,
    minimum height=0.85cm, inner sep=5pt, font=\small},
  wide/.style={box, text width=11.7cm}
]
  \node[wide] (proof) at (0,0) {Same proof, source revision, and collection budget};
  \node[box] (firsta) at (-3.25,-1.4) {First sampling A};
  \node[box] (firstb) at (3.25,-1.4) {First sampling B};
  \node[box] (poola) at (-3.25,-2.8) {Same dedup and similarity rules\\Candidate list $P_A$};
  \node[box] (poolb) at (3.25,-2.8) {Same dedup and similarity rules\\Candidate list $P_B$};
  \node[wide] (split) at (0,-4.5) {Assign shared record identities jointly to training or eval\\Keep each arm's membership; freeze one common eval set $E$};
  \node[box] (traina) at (-3.25,-6.2) {A training candidates: $P_A\setminus E$};
  \node[box] (trainb) at (3.25,-6.2) {B training candidates: $P_B\setminus E$};
  \node[box] (seconda) at (-3.25,-7.7) {Same second sampling rule and settings\\Generation list $G_A$};
  \node[box] (secondb) at (3.25,-7.7) {Same second sampling rule and settings\\Generation list $G_B$};
  \node[box] (desca) at (-3.25,-9.2) {Generate description A};
  \node[box] (descb) at (3.25,-9.2) {Generate description B};
  \node[wide] (judge) at (0,-10.8) {Evaluate both descriptions on the same frozen $E$};
  \draw[flow] (proof) -- (firsta);
  \draw[flow] (proof) -- (firstb);
  \draw[flow] (firsta) -- (poola);
  \draw[flow] (firstb) -- (poolb);
  \draw[flow] (poola) -- (split);
  \draw[flow] (poolb) -- (split);
  \draw[flow] (split) -- (traina);
  \draw[flow] (split) -- (trainb);
  \draw[flow] (traina) -- (seconda);
  \draw[flow] (trainb) -- (secondb);
  \draw[flow] (seconda) -- (desca);
  \draw[flow] (secondb) -- (descb);
  \draw[flow] (desca) -- (judge);
  \draw[flow] (descb) -- (judge);
\end{tikzpicture}
\end{center}

\paragraph{一起決定歸屬，不是混合兩邊的生成資料。}
一起查看 A、B 的候選，是為了辨認內容相同的 records，\allowbreak
統一決定它們是否屬於 evaluation。
若某筆 record 被選入共同 eval set，\allowbreak 必須從兩邊的 training 候選中排除。
不能將 B 獨有的 record 加進 A，\allowbreak 或反過來補入。
A training 與 B training 可以互相重疊，\allowbreak 但任一方都不能與 eval 重疊。

例如，下列編號代表依內容辨認的 record identity，\allowbreak 而不是檔案內的行號：
\begin{quote}
$P_A=\{r_1,r_2,r_3,r_4\}$，$P_B=\{r_2,r_3,r_5,r_6\}$。\\
若選定 $E=\{r_3,r_5\}$，則\\
A training 為 $\{r_1,r_2,r_4\}$；\\
B training 為 $\{r_2,r_6\}$。
\end{quote}
兩份 descriptions 都評 $E$，\allowbreak 即使其中的 $r_5$ 原本只存在於 B 候選池。
共同 eval 不限於 A、B 的交集，\allowbreak 但每個內容 identity 只保留一份。

\paragraph{控制最後送進 LLM 的筆數。}
在相同 tactic 內，\allowbreak 對 success 與 failure 分開決定共同的生成筆數：
\[
  n_{t,o}=\min\bigl(K_{t,o},\ |T_{A,t,o}|,\ |T_{B,t,o}|\bigr),
\]
其中 $t$ 表示 tactic，\allowbreak $o$ 表示 success 或 failure，\allowbreak
$T$ 是移除 eval 後的 training 候選，\allowbreak $K$ 是預定的生成上限。
兩邊各自用相同的 deterministic 第二次 sampling 規則選出 $n_{t,o}$ 筆。
這個設計對齊的是每個 tactic 的 success／failure 總筆數，\allowbreak
不額外要求各 goal-count group 的實際筆數都相同。
相同筆數也不代表相同 token 數或 LLM 費用。
若某個 outcome 在任一邊沒有 training 候選，\allowbreak
對齊後兩邊該 outcome 都使用零筆；\allowbreak 應明確記錄此限制。

例如，若 success 上限為 40，\allowbreak
兩邊可用候選分別為 35 與 28，\allowbreak 最後各選 28 筆。
較多的一方不必刪除其餘已保存候選；\allowbreak 它們只是未用於這輪生成。
這是\textbf{固定生成筆數的比較}，\allowbreak 不是讓每種方法用盡所有可用資料的比較。

\paragraph{第一次先保留較大的候選池。}
若擔心切出 eval 後 training 太少，\allowbreak
可以在第一次 sampling 使用相同且較大的候選容量。
例如，兩邊先各保留 120 筆候選，\allowbreak
經過去重、similarity 與 eval 排除後，\allowbreak 再各選最多 40 筆 success。
這只是示例，\allowbreak 不是已決定的參數，也不保證每個 tactic 都有足夠資料。
先離線統計各階段保留量，\allowbreak 再決定容量與 eval 大小；\allowbreak
不能複製 records 湊數或將 eval 放回 training。
候選容量增加也會增加收集與儲存成本。

\subsection{情況二：比較第二次 sampling 方法}
\label{subsec:holdout-second-sampling}

第一次 sampling、dedup 與 similarity check 全部相同，\allowbreak
先取得一份共同候選池 $P$。
\textbf{在第二次 sampling 之前只切分一次}，\allowbreak
得到共同 training 候選池 $T$ 與固定 eval set $E$。
A、B 的第二次 sampling 都從同一份 $T$ 開始選。

\begin{center}
\begin{tikzpicture}[
  >=Latex,
  flow/.style={->, thick},
  box/.style={draw, rounded corners, align=center, text width=5.2cm,
    minimum height=0.85cm, inner sep=5pt, font=\small},
  wide/.style={box, text width=11.7cm}
]
  \node[wide] (pool) at (0,0) {Same first sampling, dedup, and similarity check\\One common candidate pool $P$};
  \node[wide] (split) at (0,-1.6) {Split once, before either second sampling method};
  \node[box] (training) at (-3.25,-3.1) {Common training candidates $T$};
  \node[box] (eval) at (3.25,-3.1) {Frozen eval set $E$\\Never passed to generation};
  \node[box] (seconda) at (-3.25,-5.0) {Second sampling A on $T$\\Generation list $G_A$};
  \node[box] (secondb) at (3.25,-5.0) {Second sampling B on $T$\\Generation list $G_B$};
  \node[box] (desca) at (-3.25,-6.6) {Generate description A};
  \node[box] (descb) at (3.25,-6.6) {Generate description B};
  \node[wide] (judge) at (0,-8.3) {Evaluate both descriptions on the same frozen $E$};
  \draw[flow] (pool) -- (split);
  \draw[flow] (split) -- (training);
  \draw[flow] (split) -- (eval);
  \draw[flow] (training) -- (seconda);
  \draw[flow] (training.south) -- (-3.25,-4.05) -- (3.25,-4.05) -- (secondb.north);
  \draw[flow] (seconda) -- (desca);
  \draw[flow] (secondb) -- (descb);
  \draw[flow] (desca) -- (judge);
  \draw[flow] (descb) -- (judge);
  \draw[flow] (eval.east) -- (6.35,-3.1) -- (6.35,-8.3) -- (judge.east);
\end{tikzpicture}
\end{center}

例如，共同候選池有 100 筆 success，\allowbreak
先保留 20 筆作為 eval，\allowbreak 剩下 80 筆作為 training 候選。
A、B 各自從同一份 80 筆候選選 40 筆，\allowbreak
生成兩份 descriptions，\allowbreak 然後都評同一份預留的 20 筆。
這些數字只是示意，\allowbreak failure 也需要另外規劃保留量。

不建議先執行第二次 sampling，\allowbreak
再從 A、B 的最終 lists 抽走 eval，\allowbreak 然後又 sampling 一次。
這會移除方法刻意保留的代表，\allowbreak
測到的變成兩輪選取與移除的整體效果。
若要單獨比較第二次選取方法，\allowbreak
應固定共同 training 候選、生成上限與非比較中的設定。
若某個方法因組別限制無法填滿共同目標，\allowbreak
須事先決定如何處理並記錄實際筆數，\allowbreak 不能默認兩邊一定選滿。

\subsection{兩種設計都必須遵守的規則}

令 $G_A$、$G_B$ 為實際送進 LLM 的 generation lists，\allowbreak
共同的無重疊要求為
\[
  G_A\cap E=\emptyset,\qquad G_B\cap E=\emptyset.
\]
不要求 $G_A\cap G_B=\emptyset$：\allowbreak
兩種方法可以選中相同的 training records。

\begin{enumerate}
  \item \textbf{按內容決定歸屬。}
        不能只比較檔名、proof 名稱或 record 編號。
        應依一致的內容 identity 規則辨認 records，\allowbreak
        並檢查生成與評分實際看到的 evidence 是否重複；\allowbreak
        僅來源 metadata 不同不能當成新的獨立案例。
  \item \textbf{先切分、固定，再生成與評分。}
        切分採 deterministic 規則，\allowbreak 保存 membership 與設定，\allowbreak
        不依生成結果或 judge 分數重新挑選 eval。
        切分比例、穩定排序規則、各 tactic／outcome 最低筆數仍待決定。
  \item \textbf{不足資料時明確記錄。}
        分別記錄各 tactic 的 success／failure 與 goal-count 分布。
        若沒有可評 records，\allowbreak 該項標記 N/A；\allowbreak
        若生成資料不足，\allowbreak 依事先約定的規則縮小規模或不納入，\allowbreak
        不把缺失資料當成零分。
  \item \textbf{固定生成與 judge 條件。}
        除了要比較的 sampling 方法，\allowbreak
        固定 prompt、模型、effort、token budget 與評分規則。
        兩份 descriptions 評同一份 eval evidence，\allowbreak
        分別呈現每個 tactic 的 success／failure 分數與筆數。
        Deterministic 選取不代表 LLM 回答必然相同。
  \item \textbf{保存每階段數量與來源。}
        保存 source commit、collector／selector 版本、候選設定、\allowbreak
        去重與 similarity 設定、split membership、\allowbreak
        最終 generation lists 與固定 eval set。
        報告原本可用量及最後實際使用量，\allowbreak
        避免將筆數差異誤解為內容選取的改善。
\end{enumerate}

\paragraph{解讀限制。}
這是同一批 proof records 內的 held-out evaluation，\allowbreak
不是對另一個未見 proof 的泛化測試。
完全不同的 records 也可能非常相似；\allowbreak
若需要更嚴格的分離，\allowbreak 可將相似 records 以固定共同規則分成群，\allowbreak
整群分到同一邊，\allowbreak 但此項尚未決定。
此外，切分發生在 similarity check 之後，\allowbreak
共同 eval 的分布仍受上游收集與篩選影響，\allowbreak
不能將結果直接推廣到所有 proof 執行案例。

% Needs Chinese-capable fonts, tikz, and \usetikzlibrary{arrows.meta}.
% The existing framework.tex already provides these prerequisites.

\section{比較兩階段 sampling 時的 training 與 evaluation 切分}
\label{sec:sampling-training-evaluation-workflows}

本節記錄預定的實驗設計，\allowbreak 不代表已經實作或執行。
這裡的 training records 是提供給 LLM 生成 description 的範例，\allowbreak
不會訓練或更新模型權重。
為避免混淆，\allowbreak 下文區分「training 候選池」與「最後送進 LLM 的 generation list」。

\begin{itemize}
  \item \textbf{第一次 sampling：}collector 在 proof 執行時選取並保留候選 records。
  \item \textbf{中間處理：}依既定流程合併候選、exact dedup 與 similarity check。
  \item \textbf{第二次 sampling：}從候選池依分組配額與遞補規則，\allowbreak
        選出最後用來生成 description 的 records。
\end{itemize}

\subsection{情況一：比較第一次 sampling 方法}
\label{subsec:holdout-first-sampling}

A、B 從相同 proof 與 source commit 開始，\allowbreak
使用不同的第一次 sampling 方法。
兩邊使用相同候選容量，\allowbreak 並固定其餘 collector 設定、dedup 與 similarity check 規則。
因此，中間處理後的候選 lists 可以不同，\allowbreak 這正是要保留的實驗差異。

\begin{center}
\begin{tikzpicture}[
  >=Latex,
  flow/.style={->, thick},
  box/.style={draw, rounded corners, align=center, text width=5.2cm,
    minimum height=0.85cm, inner sep=5pt, font=\small},
  wide/.style={box, text width=11.7cm}
]
  \node[wide] (proof) at (0,0) {Same proof, source revision, and collection budget};
  \node[box] (firsta) at (-3.25,-1.4) {First sampling A};
  \node[box] (firstb) at (3.25,-1.4) {First sampling B};
  \node[box] (poola) at (-3.25,-2.8) {Same dedup and similarity rules\\Candidate list $P_A$};
  \node[box] (poolb) at (3.25,-2.8) {Same dedup and similarity rules\\Candidate list $P_B$};
  \node[wide] (split) at (0,-4.5) {Assign shared record identities jointly to training or eval\\Keep each arm's membership; freeze one common eval set $E$};
  \node[box] (traina) at (-3.25,-6.2) {A training candidates: $P_A\setminus E$};
  \node[box] (trainb) at (3.25,-6.2) {B training candidates: $P_B\setminus E$};
  \node[box] (seconda) at (-3.25,-7.7) {Same second sampling rule and settings\\Generation list $G_A$};
  \node[box] (secondb) at (3.25,-7.7) {Same second sampling rule and settings\\Generation list $G_B$};
  \node[box] (desca) at (-3.25,-9.2) {Generate description A};
  \node[box] (descb) at (3.25,-9.2) {Generate description B};
  \node[wide] (judge) at (0,-10.8) {Evaluate both descriptions on the same frozen $E$};
  \draw[flow] (proof) -- (firsta);
  \draw[flow] (proof) -- (firstb);
  \draw[flow] (firsta) -- (poola);
  \draw[flow] (firstb) -- (poolb);
  \draw[flow] (poola) -- (split);
  \draw[flow] (poolb) -- (split);
  \draw[flow] (split) -- (traina);
  \draw[flow] (split) -- (trainb);
  \draw[flow] (traina) -- (seconda);
  \draw[flow] (trainb) -- (secondb);
  \draw[flow] (seconda) -- (desca);
  \draw[flow] (secondb) -- (descb);
  \draw[flow] (desca) -- (judge);
  \draw[flow] (descb) -- (judge);
\end{tikzpicture}
\end{center}

\paragraph{一起決定歸屬，不是混合兩邊的生成資料。}
一起查看 A、B 的候選，是為了辨認內容相同的 records，\allowbreak
統一決定它們是否屬於 evaluation。
若某筆 record 被選入共同 eval set，\allowbreak 必須從兩邊的 training 候選中排除。
不能將 B 獨有的 record 加進 A，\allowbreak 或反過來補入。
A training 與 B training 可以互相重疊，\allowbreak 但任一方都不能與 eval 重疊。

例如，下列編號代表依內容辨認的 record identity，\allowbreak 而不是檔案內的行號：
\begin{quote}
$P_A=\{r_1,r_2,r_3,r_4\}$，$P_B=\{r_2,r_3,r_5,r_6\}$。\\
若選定 $E=\{r_3,r_5\}$，則\\
A training 為 $\{r_1,r_2,r_4\}$；\\
B training 為 $\{r_2,r_6\}$。
\end{quote}
兩份 descriptions 都評 $E$，\allowbreak 即使其中的 $r_5$ 原本只存在於 B 候選池。
共同 eval 不限於 A、B 的交集，\allowbreak 但每個內容 identity 只保留一份。

\paragraph{控制最後送進 LLM 的筆數。}
在相同 tactic 內，\allowbreak 對 success 與 failure 分開決定共同的生成筆數：
\[
  n_{t,o}=\min\bigl(K_{t,o},\ |T_{A,t,o}|,\ |T_{B,t,o}|\bigr),
\]
其中 $t$ 表示 tactic，\allowbreak $o$ 表示 success 或 failure，\allowbreak
$T$ 是移除 eval 後的 training 候選，\allowbreak $K$ 是預定的生成上限。
兩邊各自用相同的 deterministic 第二次 sampling 規則選出 $n_{t,o}$ 筆。
這個設計對齊的是每個 tactic 的 success／failure 總筆數，\allowbreak
不額外要求各 goal-count group 的實際筆數都相同。
相同筆數也不代表相同 token 數或 LLM 費用。
若某個 outcome 在任一邊沒有 training 候選，\allowbreak
對齊後兩邊該 outcome 都使用零筆；\allowbreak 應明確記錄此限制。

例如，若 success 上限為 40，\allowbreak
兩邊可用候選分別為 35 與 28，\allowbreak 最後各選 28 筆。
較多的一方不必刪除其餘已保存候選；\allowbreak 它們只是未用於這輪生成。
這是\textbf{固定生成筆數的比較}，\allowbreak 不是讓每種方法用盡所有可用資料的比較。

\paragraph{第一次先保留較大的候選池。}
若擔心切出 eval 後 training 太少，\allowbreak
可以在第一次 sampling 使用相同且較大的候選容量。
例如，兩邊先各保留 120 筆候選，\allowbreak
經過去重、similarity 與 eval 排除後，\allowbreak 再各選最多 40 筆 success。
這只是示例，\allowbreak 不是已決定的參數，也不保證每個 tactic 都有足夠資料。
先離線統計各階段保留量，\allowbreak 再決定容量與 eval 大小；\allowbreak
不能複製 records 湊數或將 eval 放回 training。
候選容量增加也會增加收集與儲存成本。

\subsection{情況二：比較第二次 sampling 方法}
\label{subsec:holdout-second-sampling}

第一次 sampling、dedup 與 similarity check 全部相同，\allowbreak
先取得一份共同候選池 $P$。
\textbf{在第二次 sampling 之前只切分一次}，\allowbreak
得到共同 training 候選池 $T$ 與固定 eval set $E$。
A、B 的第二次 sampling 都從同一份 $T$ 開始選。

\begin{center}
\begin{tikzpicture}[
  >=Latex,
  flow/.style={->, thick},
  box/.style={draw, rounded corners, align=center, text width=5.2cm,
    minimum height=0.85cm, inner sep=5pt, font=\small},
  wide/.style={box, text width=11.7cm}
]
  \node[wide] (pool) at (0,0) {Same first sampling, dedup, and similarity check\\One common candidate pool $P$};
  \node[wide] (split) at (0,-1.6) {Split once, before either second sampling method};
  \node[box] (training) at (-3.25,-3.1) {Common training candidates $T$};
  \node[box] (eval) at (3.25,-3.1) {Frozen eval set $E$\\Never passed to generation};
  \node[box] (seconda) at (-3.25,-5.0) {Second sampling A on $T$\\Generation list $G_A$};
  \node[box] (secondb) at (3.25,-5.0) {Second sampling B on $T$\\Generation list $G_B$};
  \node[box] (desca) at (-3.25,-6.6) {Generate description A};
  \node[box] (descb) at (3.25,-6.6) {Generate description B};
  \node[wide] (judge) at (0,-8.3) {Evaluate both descriptions on the same frozen $E$};
  \draw[flow] (pool) -- (split);
  \draw[flow] (split) -- (training);
  \draw[flow] (split) -- (eval);
  \draw[flow] (training) -- (seconda);
  \draw[flow] (training.south) -- (-3.25,-4.05) -- (3.25,-4.05) -- (secondb.north);
  \draw[flow] (seconda) -- (desca);
  \draw[flow] (secondb) -- (descb);
  \draw[flow] (desca) -- (judge);
  \draw[flow] (descb) -- (judge);
  \draw[flow] (eval.east) -- (6.35,-3.1) -- (6.35,-8.3) -- (judge.east);
\end{tikzpicture}
\end{center}

例如，共同候選池有 100 筆 success，\allowbreak
先保留 20 筆作為 eval，\allowbreak 剩下 80 筆作為 training 候選。
A、B 各自從同一份 80 筆候選選 40 筆，\allowbreak
生成兩份 descriptions，\allowbreak 然後都評同一份預留的 20 筆。
這些數字只是示意，\allowbreak failure 也需要另外規劃保留量。

不建議先執行第二次 sampling，\allowbreak
再從 A、B 的最終 lists 抽走 eval，\allowbreak 然後又 sampling 一次。
這會移除方法刻意保留的代表，\allowbreak
測到的變成兩輪選取與移除的整體效果。
若要單獨比較第二次選取方法，\allowbreak
應固定共同 training 候選、生成上限與非比較中的設定。
若某個方法因組別限制無法填滿共同目標，\allowbreak
須事先決定如何處理並記錄實際筆數，\allowbreak 不能默認兩邊一定選滿。

\subsection{兩種設計都必須遵守的規則}

令 $G_A$、$G_B$ 為實際送進 LLM 的 generation lists，\allowbreak
共同的無重疊要求為
\[
  G_A\cap E=\emptyset,\qquad G_B\cap E=\emptyset.
\]
不要求 $G_A\cap G_B=\emptyset$：\allowbreak
兩種方法可以選中相同的 training records。

\begin{enumerate}
  \item \textbf{按內容決定歸屬。}
        不能只比較檔名、proof 名稱或 record 編號。
        應依一致的內容 identity 規則辨認 records，\allowbreak
        並檢查生成與評分實際看到的 evidence 是否重複；\allowbreak
        僅來源 metadata 不同不能當成新的獨立案例。
  \item \textbf{先切分、固定，再生成與評分。}
        切分採 deterministic 規則，\allowbreak 保存 membership 與設定，\allowbreak
        不依生成結果或 judge 分數重新挑選 eval。
        切分比例、穩定排序規則、各 tactic／outcome 最低筆數仍待決定。
  \item \textbf{不足資料時明確記錄。}
        分別記錄各 tactic 的 success／failure 與 goal-count 分布。
        若沒有可評 records，\allowbreak 該項標記 N/A；\allowbreak
        若生成資料不足，\allowbreak 依事先約定的規則縮小規模或不納入，\allowbreak
        不把缺失資料當成零分。
  \item \textbf{固定生成與 judge 條件。}
        除了要比較的 sampling 方法，\allowbreak
        固定 prompt、模型、effort、token budget 與評分規則。
        兩份 descriptions 評同一份 eval evidence，\allowbreak
        分別呈現每個 tactic 的 success／failure 分數與筆數。
        Deterministic 選取不代表 LLM 回答必然相同。
  \item \textbf{保存每階段數量與來源。}
        保存 source commit、collector／selector 版本、候選設定、\allowbreak
        去重與 similarity 設定、split membership、\allowbreak
        最終 generation lists 與固定 eval set。
        報告原本可用量及最後實際使用量，\allowbreak
        避免將筆數差異誤解為內容選取的改善。
\end{enumerate}

\paragraph{解讀限制。}
這是同一批 proof records 內的 held-out evaluation，\allowbreak
不是對另一個未見 proof 的泛化測試。
完全不同的 records 也可能非常相似；\allowbreak
若需要更嚴格的分離，\allowbreak 可將相似 records 以固定共同規則分成群，\allowbreak
整群分到同一邊，\allowbreak 但此項尚未決定。
此外，切分發生在 similarity check 之後，\allowbreak
共同 eval 的分布仍受上游收集與篩選影響，\allowbreak
不能將結果直接推廣到所有 proof 執行案例。

% Needs Chinese-capable fonts, tikz, and \usetikzlibrary{arrows.meta}.
% The existing framework.tex already provides these prerequisites.

\section{比較兩階段 sampling 時的 training 與 evaluation 切分}
\label{sec:sampling-training-evaluation-workflows}

本節記錄預定的實驗設計，\allowbreak 不代表已經實作或執行。
這裡的 training records 是提供給 LLM 生成 description 的範例，\allowbreak
不會訓練或更新模型權重。
為避免混淆，\allowbreak 下文區分「training 候選池」與「最後送進 LLM 的 generation list」。

\begin{itemize}
  \item \textbf{第一次 sampling：}collector 在 proof 執行時選取並保留候選 records。
  \item \textbf{中間處理：}依既定流程合併候選、exact dedup 與 similarity check。
  \item \textbf{第二次 sampling：}從候選池依分組配額與遞補規則，\allowbreak
        選出最後用來生成 description 的 records。
\end{itemize}

\subsection{情況一：比較第一次 sampling 方法}
\label{subsec:holdout-first-sampling}

A、B 從相同 proof 與 source commit 開始，\allowbreak
使用不同的第一次 sampling 方法。
兩邊使用相同候選容量，\allowbreak 並固定其餘 collector 設定、dedup 與 similarity check 規則。
因此，中間處理後的候選 lists 可以不同，\allowbreak 這正是要保留的實驗差異。

\begin{center}
\begin{tikzpicture}[
  >=Latex,
  flow/.style={->, thick},
  box/.style={draw, rounded corners, align=center, text width=5.2cm,
    minimum height=0.85cm, inner sep=5pt, font=\small},
  wide/.style={box, text width=11.7cm}
]
  \node[wide] (proof) at (0,0) {Same proof, source revision, and collection budget};
  \node[box] (firsta) at (-3.25,-1.4) {First sampling A};
  \node[box] (firstb) at (3.25,-1.4) {First sampling B};
  \node[box] (poola) at (-3.25,-2.8) {Same dedup and similarity rules\\Candidate list $P_A$};
  \node[box] (poolb) at (3.25,-2.8) {Same dedup and similarity rules\\Candidate list $P_B$};
  \node[wide] (split) at (0,-4.5) {Assign shared record identities jointly to training or eval\\Keep each arm's membership; freeze one common eval set $E$};
  \node[box] (traina) at (-3.25,-6.2) {A training candidates: $P_A\setminus E$};
  \node[box] (trainb) at (3.25,-6.2) {B training candidates: $P_B\setminus E$};
  \node[box] (seconda) at (-3.25,-7.7) {Same second sampling rule and settings\\Generation list $G_A$};
  \node[box] (secondb) at (3.25,-7.7) {Same second sampling rule and settings\\Generation list $G_B$};
  \node[box] (desca) at (-3.25,-9.2) {Generate description A};
  \node[box] (descb) at (3.25,-9.2) {Generate description B};
  \node[wide] (judge) at (0,-10.8) {Evaluate both descriptions on the same frozen $E$};
  \draw[flow] (proof) -- (firsta);
  \draw[flow] (proof) -- (firstb);
  \draw[flow] (firsta) -- (poola);
  \draw[flow] (firstb) -- (poolb);
  \draw[flow] (poola) -- (split);
  \draw[flow] (poolb) -- (split);
  \draw[flow] (split) -- (traina);
  \draw[flow] (split) -- (trainb);
  \draw[flow] (traina) -- (seconda);
  \draw[flow] (trainb) -- (secondb);
  \draw[flow] (seconda) -- (desca);
  \draw[flow] (secondb) -- (descb);
  \draw[flow] (desca) -- (judge);
  \draw[flow] (descb) -- (judge);
\end{tikzpicture}
\end{center}

\paragraph{一起決定歸屬，不是混合兩邊的生成資料。}
一起查看 A、B 的候選，是為了辨認內容相同的 records，\allowbreak
統一決定它們是否屬於 evaluation。
若某筆 record 被選入共同 eval set，\allowbreak 必須從兩邊的 training 候選中排除。
不能將 B 獨有的 record 加進 A，\allowbreak 或反過來補入。
A training 與 B training 可以互相重疊，\allowbreak 但任一方都不能與 eval 重疊。

例如，下列編號代表依內容辨認的 record identity，\allowbreak 而不是檔案內的行號：
\begin{quote}
$P_A=\{r_1,r_2,r_3,r_4\}$，$P_B=\{r_2,r_3,r_5,r_6\}$。\\
若選定 $E=\{r_3,r_5\}$，則\\
A training 為 $\{r_1,r_2,r_4\}$；\\
B training 為 $\{r_2,r_6\}$。
\end{quote}
兩份 descriptions 都評 $E$，\allowbreak 即使其中的 $r_5$ 原本只存在於 B 候選池。
共同 eval 不限於 A、B 的交集，\allowbreak 但每個內容 identity 只保留一份。

\paragraph{控制最後送進 LLM 的筆數。}
在相同 tactic 內，\allowbreak 對 success 與 failure 分開決定共同的生成筆數：
\[
  n_{t,o}=\min\bigl(K_{t,o},\ |T_{A,t,o}|,\ |T_{B,t,o}|\bigr),
\]
其中 $t$ 表示 tactic，\allowbreak $o$ 表示 success 或 failure，\allowbreak
$T$ 是移除 eval 後的 training 候選，\allowbreak $K$ 是預定的生成上限。
兩邊各自用相同的 deterministic 第二次 sampling 規則選出 $n_{t,o}$ 筆。
這個設計對齊的是每個 tactic 的 success／failure 總筆數，\allowbreak
不額外要求各 goal-count group 的實際筆數都相同。
相同筆數也不代表相同 token 數或 LLM 費用。
若某個 outcome 在任一邊沒有 training 候選，\allowbreak
對齊後兩邊該 outcome 都使用零筆；\allowbreak 應明確記錄此限制。

例如，若 success 上限為 40，\allowbreak
兩邊可用候選分別為 35 與 28，\allowbreak 最後各選 28 筆。
較多的一方不必刪除其餘已保存候選；\allowbreak 它們只是未用於這輪生成。
這是\textbf{固定生成筆數的比較}，\allowbreak 不是讓每種方法用盡所有可用資料的比較。

\paragraph{第一次先保留較大的候選池。}
若擔心切出 eval 後 training 太少，\allowbreak
可以在第一次 sampling 使用相同且較大的候選容量。
例如，兩邊先各保留 120 筆候選，\allowbreak
經過去重、similarity 與 eval 排除後，\allowbreak 再各選最多 40 筆 success。
這只是示例，\allowbreak 不是已決定的參數，也不保證每個 tactic 都有足夠資料。
先離線統計各階段保留量，\allowbreak 再決定容量與 eval 大小；\allowbreak
不能複製 records 湊數或將 eval 放回 training。
候選容量增加也會增加收集與儲存成本。

\subsection{情況二：比較第二次 sampling 方法}
\label{subsec:holdout-second-sampling}

第一次 sampling、dedup 與 similarity check 全部相同，\allowbreak
先取得一份共同候選池 $P$。
\textbf{在第二次 sampling 之前只切分一次}，\allowbreak
得到共同 training 候選池 $T$ 與固定 eval set $E$。
A、B 的第二次 sampling 都從同一份 $T$ 開始選。

\begin{center}
\begin{tikzpicture}[
  >=Latex,
  flow/.style={->, thick},
  box/.style={draw, rounded corners, align=center, text width=5.2cm,
    minimum height=0.85cm, inner sep=5pt, font=\small},
  wide/.style={box, text width=11.7cm}
]
  \node[wide] (pool) at (0,0) {Same first sampling, dedup, and similarity check\\One common candidate pool $P$};
  \node[wide] (split) at (0,-1.6) {Split once, before either second sampling method};
  \node[box] (training) at (-3.25,-3.1) {Common training candidates $T$};
  \node[box] (eval) at (3.25,-3.1) {Frozen eval set $E$\\Never passed to generation};
  \node[box] (seconda) at (-3.25,-5.0) {Second sampling A on $T$\\Generation list $G_A$};
  \node[box] (secondb) at (3.25,-5.0) {Second sampling B on $T$\\Generation list $G_B$};
  \node[box] (desca) at (-3.25,-6.6) {Generate description A};
  \node[box] (descb) at (3.25,-6.6) {Generate description B};
  \node[wide] (judge) at (0,-8.3) {Evaluate both descriptions on the same frozen $E$};
  \draw[flow] (pool) -- (split);
  \draw[flow] (split) -- (training);
  \draw[flow] (split) -- (eval);
  \draw[flow] (training) -- (seconda);
  \draw[flow] (training.south) -- (-3.25,-4.05) -- (3.25,-4.05) -- (secondb.north);
  \draw[flow] (seconda) -- (desca);
  \draw[flow] (secondb) -- (descb);
  \draw[flow] (desca) -- (judge);
  \draw[flow] (descb) -- (judge);
  \draw[flow] (eval.east) -- (6.35,-3.1) -- (6.35,-8.3) -- (judge.east);
\end{tikzpicture}
\end{center}

例如，共同候選池有 100 筆 success，\allowbreak
先保留 20 筆作為 eval，\allowbreak 剩下 80 筆作為 training 候選。
A、B 各自從同一份 80 筆候選選 40 筆，\allowbreak
生成兩份 descriptions，\allowbreak 然後都評同一份預留的 20 筆。
這些數字只是示意，\allowbreak failure 也需要另外規劃保留量。

不建議先執行第二次 sampling，\allowbreak
再從 A、B 的最終 lists 抽走 eval，\allowbreak 然後又 sampling 一次。
這會移除方法刻意保留的代表，\allowbreak
測到的變成兩輪選取與移除的整體效果。
若要單獨比較第二次選取方法，\allowbreak
應固定共同 training 候選、生成上限與非比較中的設定。
若某個方法因組別限制無法填滿共同目標，\allowbreak
須事先決定如何處理並記錄實際筆數，\allowbreak 不能默認兩邊一定選滿。

\subsection{兩種設計都必須遵守的規則}

令 $G_A$、$G_B$ 為實際送進 LLM 的 generation lists，\allowbreak
共同的無重疊要求為
\[
  G_A\cap E=\emptyset,\qquad G_B\cap E=\emptyset.
\]
不要求 $G_A\cap G_B=\emptyset$：\allowbreak
兩種方法可以選中相同的 training records。

\begin{enumerate}
  \item \textbf{按內容決定歸屬。}
        不能只比較檔名、proof 名稱或 record 編號。
        應依一致的內容 identity 規則辨認 records，\allowbreak
        並檢查生成與評分實際看到的 evidence 是否重複；\allowbreak
        僅來源 metadata 不同不能當成新的獨立案例。
  \item \textbf{先切分、固定，再生成與評分。}
        切分採 deterministic 規則，\allowbreak 保存 membership 與設定，\allowbreak
        不依生成結果或 judge 分數重新挑選 eval。
        切分比例、穩定排序規則、各 tactic／outcome 最低筆數仍待決定。
  \item \textbf{不足資料時明確記錄。}
        分別記錄各 tactic 的 success／failure 與 goal-count 分布。
        若沒有可評 records，\allowbreak 該項標記 N/A；\allowbreak
        若生成資料不足，\allowbreak 依事先約定的規則縮小規模或不納入，\allowbreak
        不把缺失資料當成零分。
  \item \textbf{固定生成與 judge 條件。}
        除了要比較的 sampling 方法，\allowbreak
        固定 prompt、模型、effort、token budget 與評分規則。
        兩份 descriptions 評同一份 eval evidence，\allowbreak
        分別呈現每個 tactic 的 success／failure 分數與筆數。
        Deterministic 選取不代表 LLM 回答必然相同。
  \item \textbf{保存每階段數量與來源。}
        保存 source commit、collector／selector 版本、候選設定、\allowbreak
        去重與 similarity 設定、split membership、\allowbreak
        最終 generation lists 與固定 eval set。
        報告原本可用量及最後實際使用量，\allowbreak
        避免將筆數差異誤解為內容選取的改善。
\end{enumerate}

\paragraph{解讀限制。}
這是同一批 proof records 內的 held-out evaluation，\allowbreak
不是對另一個未見 proof 的泛化測試。
完全不同的 records 也可能非常相似；\allowbreak
若需要更嚴格的分離，\allowbreak 可將相似 records 以固定共同規則分成群，\allowbreak
整群分到同一邊，\allowbreak 但此項尚未決定。
此外，切分發生在 similarity check 之後，\allowbreak
共同 eval 的分布仍受上游收集與篩選影響，\allowbreak
不能將結果直接推廣到所有 proof 執行案例。

% Needs Chinese-capable fonts, tikz, and \usetikzlibrary{arrows.meta}.
% The existing framework.tex already provides these prerequisites.

\section{比較兩階段 sampling 時的 training 與 evaluation 切分}
\label{sec:sampling-training-evaluation-workflows}

本節記錄預定的實驗設計，\allowbreak 不代表已經實作或執行。
這裡的 training records 是提供給 LLM 生成 description 的範例，\allowbreak
不會訓練或更新模型權重。
為避免混淆，\allowbreak 下文區分「training 候選池」與「最後送進 LLM 的 generation list」。

\begin{itemize}
  \item \textbf{第一次 sampling：}collector 在 proof 執行時選取並保留候選 records。
  \item \textbf{中間處理：}依既定流程合併候選、exact dedup 與 similarity check。
  \item \textbf{第二次 sampling：}從候選池依分組配額與遞補規則，\allowbreak
        選出最後用來生成 description 的 records。
\end{itemize}

\subsection{情況一：比較第一次 sampling 方法}
\label{subsec:holdout-first-sampling}

A、B 從相同 proof 與 source commit 開始，\allowbreak
使用不同的第一次 sampling 方法。
兩邊使用相同候選容量，\allowbreak 並固定其餘 collector 設定、dedup 與 similarity check 規則。
因此，中間處理後的候選 lists 可以不同，\allowbreak 這正是要保留的實驗差異。

\begin{center}
\begin{tikzpicture}[
  >=Latex,
  flow/.style={->, thick},
  box/.style={draw, rounded corners, align=center, text width=5.2cm,
    minimum height=0.85cm, inner sep=5pt, font=\small},
  wide/.style={box, text width=11.7cm}
]
  \node[wide] (proof) at (0,0) {Same proof, source revision, and collection budget};
  \node[box] (firsta) at (-3.25,-1.4) {First sampling A};
  \node[box] (firstb) at (3.25,-1.4) {First sampling B};
  \node[box] (poola) at (-3.25,-2.8) {Same dedup and similarity rules\\Candidate list $P_A$};
  \node[box] (poolb) at (3.25,-2.8) {Same dedup and similarity rules\\Candidate list $P_B$};
  \node[wide] (split) at (0,-4.5) {Assign shared record identities jointly to training or eval\\Keep each arm's membership; freeze one common eval set $E$};
  \node[box] (traina) at (-3.25,-6.2) {A training candidates: $P_A\setminus E$};
  \node[box] (trainb) at (3.25,-6.2) {B training candidates: $P_B\setminus E$};
  \node[box] (seconda) at (-3.25,-7.7) {Same second sampling rule and settings\\Generation list $G_A$};
  \node[box] (secondb) at (3.25,-7.7) {Same second sampling rule and settings\\Generation list $G_B$};
  \node[box] (desca) at (-3.25,-9.2) {Generate description A};
  \node[box] (descb) at (3.25,-9.2) {Generate description B};
  \node[wide] (judge) at (0,-10.8) {Evaluate both descriptions on the same frozen $E$};
  \draw[flow] (proof) -- (firsta);
  \draw[flow] (proof) -- (firstb);
  \draw[flow] (firsta) -- (poola);
  \draw[flow] (firstb) -- (poolb);
  \draw[flow] (poola) -- (split);
  \draw[flow] (poolb) -- (split);
  \draw[flow] (split) -- (traina);
  \draw[flow] (split) -- (trainb);
  \draw[flow] (traina) -- (seconda);
  \draw[flow] (trainb) -- (secondb);
  \draw[flow] (seconda) -- (desca);
  \draw[flow] (secondb) -- (descb);
  \draw[flow] (desca) -- (judge);
  \draw[flow] (descb) -- (judge);
\end{tikzpicture}
\end{center}

\paragraph{一起決定歸屬，不是混合兩邊的生成資料。}
一起查看 A、B 的候選，是為了辨認內容相同的 records，\allowbreak
統一決定它們是否屬於 evaluation。
若某筆 record 被選入共同 eval set，\allowbreak 必須從兩邊的 training 候選中排除。
不能將 B 獨有的 record 加進 A，\allowbreak 或反過來補入。
A training 與 B training 可以互相重疊，\allowbreak 但任一方都不能與 eval 重疊。

例如，下列編號代表依內容辨認的 record identity，\allowbreak 而不是檔案內的行號：
\begin{quote}
$P_A=\{r_1,r_2,r_3,r_4\}$，$P_B=\{r_2,r_3,r_5,r_6\}$。\\
若選定 $E=\{r_3,r_5\}$，則\\
A training 為 $\{r_1,r_2,r_4\}$；\\
B training 為 $\{r_2,r_6\}$。
\end{quote}
兩份 descriptions 都評 $E$，\allowbreak 即使其中的 $r_5$ 原本只存在於 B 候選池。
共同 eval 不限於 A、B 的交集，\allowbreak 但每個內容 identity 只保留一份。

\paragraph{控制最後送進 LLM 的筆數。}
在相同 tactic 內，\allowbreak 對 success 與 failure 分開決定共同的生成筆數：
\[
  n_{t,o}=\min\bigl(K_{t,o},\ |T_{A,t,o}|,\ |T_{B,t,o}|\bigr),
\]
其中 $t$ 表示 tactic，\allowbreak $o$ 表示 success 或 failure，\allowbreak
$T$ 是移除 eval 後的 training 候選，\allowbreak $K$ 是預定的生成上限。
兩邊各自用相同的 deterministic 第二次 sampling 規則選出 $n_{t,o}$ 筆。
這個設計對齊的是每個 tactic 的 success／failure 總筆數，\allowbreak
不額外要求各 goal-count group 的實際筆數都相同。
相同筆數也不代表相同 token 數或 LLM 費用。
若某個 outcome 在任一邊沒有 training 候選，\allowbreak
對齊後兩邊該 outcome 都使用零筆；\allowbreak 應明確記錄此限制。

例如，若 success 上限為 40，\allowbreak
兩邊可用候選分別為 35 與 28，\allowbreak 最後各選 28 筆。
較多的一方不必刪除其餘已保存候選；\allowbreak 它們只是未用於這輪生成。
這是\textbf{固定生成筆數的比較}，\allowbreak 不是讓每種方法用盡所有可用資料的比較。

\paragraph{第一次先保留較大的候選池。}
若擔心切出 eval 後 training 太少，\allowbreak
可以在第一次 sampling 使用相同且較大的候選容量。
例如，兩邊先各保留 120 筆候選，\allowbreak
經過去重、similarity 與 eval 排除後，\allowbreak 再各選最多 40 筆 success。
這只是示例，\allowbreak 不是已決定的參數，也不保證每個 tactic 都有足夠資料。
先離線統計各階段保留量，\allowbreak 再決定容量與 eval 大小；\allowbreak
不能複製 records 湊數或將 eval 放回 training。
候選容量增加也會增加收集與儲存成本。

\subsection{情況二：比較第二次 sampling 方法}
\label{subsec:holdout-second-sampling}

第一次 sampling、dedup 與 similarity check 全部相同，\allowbreak
先取得一份共同候選池 $P$。
\textbf{在第二次 sampling 之前只切分一次}，\allowbreak
得到共同 training 候選池 $T$ 與固定 eval set $E$。
A、B 的第二次 sampling 都從同一份 $T$ 開始選。

\begin{center}
\begin{tikzpicture}[
  >=Latex,
  flow/.style={->, thick},
  box/.style={draw, rounded corners, align=center, text width=5.2cm,
    minimum height=0.85cm, inner sep=5pt, font=\small},
  wide/.style={box, text width=11.7cm}
]
  \node[wide] (pool) at (0,0) {Same first sampling, dedup, and similarity check\\One common candidate pool $P$};
  \node[wide] (split) at (0,-1.6) {Split once, before either second sampling method};
  \node[box] (training) at (-3.25,-3.1) {Common training candidates $T$};
  \node[box] (eval) at (3.25,-3.1) {Frozen eval set $E$\\Never passed to generation};
  \node[box] (seconda) at (-3.25,-5.0) {Second sampling A on $T$\\Generation list $G_A$};
  \node[box] (secondb) at (3.25,-5.0) {Second sampling B on $T$\\Generation list $G_B$};
  \node[box] (desca) at (-3.25,-6.6) {Generate description A};
  \node[box] (descb) at (3.25,-6.6) {Generate description B};
  \node[wide] (judge) at (0,-8.3) {Evaluate both descriptions on the same frozen $E$};
  \draw[flow] (pool) -- (split);
  \draw[flow] (split) -- (training);
  \draw[flow] (split) -- (eval);
  \draw[flow] (training) -- (seconda);
  \draw[flow] (training.south) -- (-3.25,-4.05) -- (3.25,-4.05) -- (secondb.north);
  \draw[flow] (seconda) -- (desca);
  \draw[flow] (secondb) -- (descb);
  \draw[flow] (desca) -- (judge);
  \draw[flow] (descb) -- (judge);
  \draw[flow] (eval.east) -- (6.35,-3.1) -- (6.35,-8.3) -- (judge.east);
\end{tikzpicture}
\end{center}

例如，共同候選池有 100 筆 success，\allowbreak
先保留 20 筆作為 eval，\allowbreak 剩下 80 筆作為 training 候選。
A、B 各自從同一份 80 筆候選選 40 筆，\allowbreak
生成兩份 descriptions，\allowbreak 然後都評同一份預留的 20 筆。
這些數字只是示意，\allowbreak failure 也需要另外規劃保留量。

不建議先執行第二次 sampling，\allowbreak
再從 A、B 的最終 lists 抽走 eval，\allowbreak 然後又 sampling 一次。
這會移除方法刻意保留的代表，\allowbreak
測到的變成兩輪選取與移除的整體效果。
若要單獨比較第二次選取方法，\allowbreak
應固定共同 training 候選、生成上限與非比較中的設定。
若某個方法因組別限制無法填滿共同目標，\allowbreak
須事先決定如何處理並記錄實際筆數，\allowbreak 不能默認兩邊一定選滿。

\subsection{兩種設計都必須遵守的規則}

令 $G_A$、$G_B$ 為實際送進 LLM 的 generation lists，\allowbreak
共同的無重疊要求為
\[
  G_A\cap E=\emptyset,\qquad G_B\cap E=\emptyset.
\]
不要求 $G_A\cap G_B=\emptyset$：\allowbreak
兩種方法可以選中相同的 training records。

\begin{enumerate}
  \item \textbf{按內容決定歸屬。}
        不能只比較檔名、proof 名稱或 record 編號。
        應依一致的內容 identity 規則辨認 records，\allowbreak
        並檢查生成與評分實際看到的 evidence 是否重複；\allowbreak
        僅來源 metadata 不同不能當成新的獨立案例。
  \item \textbf{先切分、固定，再生成與評分。}
        切分採 deterministic 規則，\allowbreak 保存 membership 與設定，\allowbreak
        不依生成結果或 judge 分數重新挑選 eval。
        切分比例、穩定排序規則、各 tactic／outcome 最低筆數仍待決定。
  \item \textbf{不足資料時明確記錄。}
        分別記錄各 tactic 的 success／failure 與 goal-count 分布。
        若沒有可評 records，\allowbreak 該項標記 N/A；\allowbreak
        若生成資料不足，\allowbreak 依事先約定的規則縮小規模或不納入，\allowbreak
        不把缺失資料當成零分。
  \item \textbf{固定生成與 judge 條件。}
        除了要比較的 sampling 方法，\allowbreak
        固定 prompt、模型、effort、token budget 與評分規則。
        兩份 descriptions 評同一份 eval evidence，\allowbreak
        分別呈現每個 tactic 的 success／failure 分數與筆數。
        Deterministic 選取不代表 LLM 回答必然相同。
  \item \textbf{保存每階段數量與來源。}
        保存 source commit、collector／selector 版本、候選設定、\allowbreak
        去重與 similarity 設定、split membership、\allowbreak
        最終 generation lists 與固定 eval set。
        報告原本可用量及最後實際使用量，\allowbreak
        避免將筆數差異誤解為內容選取的改善。
\end{enumerate}

\paragraph{解讀限制。}
這是同一批 proof records 內的 held-out evaluation，\allowbreak
不是對另一個未見 proof 的泛化測試。
完全不同的 records 也可能非常相似；\allowbreak
若需要更嚴格的分離，\allowbreak 可將相似 records 以固定共同規則分成群，\allowbreak
整群分到同一邊，\allowbreak 但此項尚未決定。
此外，切分發生在 similarity check 之後，\allowbreak
共同 eval 的分布仍受上游收集與篩選影響，\allowbreak
不能將結果直接推廣到所有 proof 執行案例。
